% Méthode et conventions de lecture
\section{Méthode et conventions}
\label{sec:methode}

\paragraph{Corpus.} Ce document ne s'appuie que sur la base \texttt{drone} de
PaperDatabase telle qu'elle était le 1\textsuperscript{er} octobre 2026 : 86 articles,
tous dotés d'une fiche de lecture complète, et 352 lignes de mesures dans
\texttt{benchmarks.csv}, issues de 67 articles. La première version (29 septembre)
reposait sur 73 articles et 324 lignes ; les 13 articles ajoutés portent surtout sur la
navigation, le \gl{SoC}{SoC} \glslink{gap9}{GAP9} et les modèles vision-langage embarqués. Un 87\textsuperscript{e}
article de la base, consacré au guidage terminal d'un drone vers une cible, n'est pas
traité. Aucune référence extérieure à la base n'est citée comme source. Les
travaux cités par un survey sont signalés comme tels ; leurs chiffres ne sont pas vérifiés
à la source et servent à orienter la lecture, pas à comparer.

\begin{table}[!htbp]
\centering\footnotesize
\caption{Composition du corpus. Un article peut relever de plusieurs axes.}
\label{tab:corpus}
\begin{tabularx}{\textwidth}{@{}LrL r@{}}
\toprule
Axe & Articles & Axe & Articles \\
\midrule
Détection drone, petits objets & 29 & Nano-drones, \gl{MCU}{MCU}, \gl{ASIC}{ASIC}, neuromorphique & 20 \\
Suivi et navigation & 37 & Accélération \gl{FPGA}{FPGA} / \gl{SoC}{SoC} & 14 \\
Contre-drone & 19 & Co-design calcul--énergie--vol & 19 \\
Jeux de données et challenges & 16 & Surveys & 8 \\
 & & Fiabilité du calcul embarqué & 2 \\
\midrule
\multicolumn{4}{@{}l}{Types : 30 articles de journal, 7 de conférence, 49 \gl{preprint}{preprints} arXiv
(dont plusieurs publiés depuis, voir le texte).} \\
\multicolumn{4}{@{}l}{Années : 2016--2020 : 14 ; 2021--2023 : 38 ; 2024--2026 : 34.} \\
\bottomrule
\end{tabularx}
\end{table}

\paragraph{Origine des chiffres.} Chaque valeur est reprise de la fiche ou de
\texttt{benchmarks.csv}, avec la section ou le tableau de l'article entre crochets, par
exemple \src{Table~4}. Une valeur que les auteurs ne donnent pas s'écrit \nr{} : elle
n'est jamais estimée. Un calcul fait à partir des chiffres publiés (par exemple un \glslink{fps}{FPS}/W)
est marqué \deduit. Sur les 352 lignes, 177 sont en confiance \emph{high} et 175 en
\emph{medium} : valeurs lues dans le texte qui commente une figure, ou tableaux
reconstruits automatiquement. Ces dernières sont marquées \med{} dans les tableaux. Aucune
ligne n'est en confiance \emph{low}.

\paragraph{Portée de la puissance.} Une puissance n'a de sens qu'avec sa \glslink{portee}{portée}. La base
distingue \emph{board} (carte entière, au wattmètre), \emph{chip} (rails de la puce ou
\gl{telemetrie}{télémétrie} embarquée comme \texttt{tegrastats}), \emph{estimated} (outil de conception,
par exemple Vivado) et \emph{unspecified} (méthode non décrite). Sur les 352 lignes,
247 n'ont pas de puissance (70\,\%) ; les autres se répartissent en 51 \emph{board},
20 \emph{chip}, 32 \emph{unspecified} et 2 \emph{estimated}. Aucune des 28 lignes
ajoutées le 1\textsuperscript{er} octobre n'a de puissance mesurée à la carte. Deux puissances de portées
différentes ne se comparent pas.

\paragraph{Gamme de puissance.} Tout résultat est situé dans une gamme : \emph{nano}
(\glslink{gap8}{GAP8}, MCU, ASIC ; de quelques \glslink{u-mw}{mW} à quelques centaines de mW), \emph{FPGA} (SoC \glslink{zynq}{Zynq} ou
Kria, de 1 à 30\,\glslink{u-w}{W} selon la carte), \emph{\glslink{jetson}{Jetson}} (autour de 10\,W, de 5 à 30\,W selon le
module et le mode). Les résultats sur \gl{GPU}{GPU} de bureau (RTX, Titan, Tesla) sont signalés
comme tels : ils ne sont jamais des résultats embarqués, quoi qu'en dise leur titre.

\paragraph{Métriques de précision.} Elles ne se comparent qu'à protocole égal. L'\glslink{iou}{IoU} moyen
mono-objet du \glslink{dac-sdc}{DAC-SDC} n'est pas une \glslink{map}{mAP} ; la \glslink{map50}{mAP@0.5} de VisDrone n'est pas l'\glslink{ap}{AP} \gl{COCO}{COCO}
[.5:.95] du protocole officiel ; l'\glslink{ap}{AP} d'UAVDT est calculée à \glslink{iou}{IoU} 0,7 dans l'article
d'origine ; certains travaux de contre-drone évaluent à \glslink{iou}{IoU} 0,1 ou 0,25.

\paragraph{Contre-drone : point de vue du capteur.} Chaque article de contre-drone est
classé \emph{air $\rightarrow$ air} (détecteur à bord d'un drone), \emph{sol
$\rightarrow$ air} (capteur fixe au sol) ou point de vue non déclaré. Seul le premier
relève de l'\gl{AI}{IA} embarquée ; les autres sont cités comme jeux de données, métriques ou
contrastes.

\paragraph{Lu contre déduit.} Sauf mention, une affirmation résume ce que disent les
auteurs. Une inférence de la fiche ou de ce document est marquée \deduit. Un article
marqué \pp{} est un preprint non relu par les pairs dans la version ingérée ; quand une
version publiée existe probablement, le texte le dit.

\paragraph{Limites de la méthode.} Le corpus est petit et vient d'une bibliographie de
départ constituée à partir de résumés. Il compte peu d'articles FPGA propres au drone
(10 papiers mesurés sur FPGA) et seulement deux sur la fiabilité. Un \enquote{gap}
signifie qu'un sujet est absent \emph{de ce corpus} ; son absence de la littérature reste
une hypothèse tant qu'une recherche ciblée ne l'a pas confirmée. Les fiches ont été
rédigées à partir du texte extrait des PDF : les valeurs données seulement dans des
figures sont souvent \nr.
